\documentclass[a4paper,14pt]{extarticle}

\usepackage{fontspec}
\usepackage{polyglossia}
\setdefaultlanguage{russian}
\setmainfont{Times New Roman}
\setmonofont{Consolas}

\usepackage{geometry}
\usepackage{xcolor}
\usepackage{graphicx}
\usepackage{float}
\usepackage{listings}
\usepackage{enumitem}
\usepackage{hyperref}

\geometry{left=25mm,right=15mm,top=20mm,bottom=20mm}
\setlength{\parindent}{1.25cm}
\setlength{\parskip}{0pt}
\sloppy
\hypersetup{colorlinks=true,linkcolor=black,urlcolor=blue}
\graphicspath{{report/screenshots/}{screenshots/}}

\lstdefinestyle{java}{
    language=Java,
    basicstyle=\ttfamily\scriptsize,
    keywordstyle=\color{blue!75!black},
    stringstyle=\color{green!45!black},
    commentstyle=\color{gray},
    showstringspaces=false,
    breaklines=true,
    breakatwhitespace=false,
    tabsize=4,
    frame=none,
    columns=fullflexible
}

\newcommand{\reportfigure}[2]{%
    \begin{figure}[H]
        \centering
        \IfFileExists{report/screenshots/#1}{%
            \includegraphics[width=0.94\linewidth,height=0.78\textheight,keepaspectratio]{report/screenshots/#1}%
        }{%
            \IfFileExists{screenshots/#1}{%
                \includegraphics[width=0.94\linewidth,height=0.78\textheight,keepaspectratio]{screenshots/#1}%
            }{%
                \fbox{\begin{minipage}[c][45mm][c]{0.88\linewidth}
                    \centering
                    \textbf{Место для реального скриншота}\\[2mm]
                    \texttt{report/screenshots/#1}
                \end{minipage}}%
            }%
        }%
        \caption{#2}
    \end{figure}
}

\begin{document}

\begin{titlepage}
    \thispagestyle{empty}
    \begin{center}
        {\color{gray!45}\small
        Федеральное государственное автономное образовательное учреждение высшего образования\\
        «Национальный исследовательский университет ИТМО»\par}

        \vspace{0.9cm}
        Факультет программной инженерии и компьютерной техники

        \vspace{0.9cm}
        Основы программной инженерии

        \vfill
        {\Large Лабораторная работа №4\par}
        \vspace{0.7cm}
        {\large Вариант 666999\par}
    \end{center}

    \vfill
    \begin{flushright}
        Выполнил:\\[5mm]
        \textit{Евграфов Артём Андреевич}\\[5mm]
        Группа P3209\\[7mm]
        Проверил:\\[5mm]
        Преподаватель основ программной инженерии\\[3mm]
        \textit{Ермаков Михаил Константинович}
    \end{flushright}

    \vfill
    \begin{center}
        {\color{gray!45}Санкт-Петербург 2026}
    \end{center}
\end{titlepage}

\tableofcontents
\newpage

\section*{Текст задания}
\addcontentsline{toc}{section}{Текст задания}

Для программы из лабораторной работы №3 по дисциплине «Веб-программирование» необходимо выполнить следующие задания.

\begin{enumerate}[leftmargin=1.2cm]
    \item Реализовать:
    \begin{itemize}
        \item MBean, считающий общее число установленных пользователем точек и число точек, не попадающих в область. Если количество установленных пользователем точек стало кратно 10, MBean должен отправить оповещение об этом событии;
        \item MBean, определяющий площадь получившейся фигуры.
    \end{itemize}

    \item С помощью утилиты JConsole провести мониторинг программы:
    \begin{itemize}
        \item снять показания MBean-классов, разработанных при выполнении задания 1;
        \item определить время в миллисекундах, прошедшее с момента запуска виртуальной машины.
    \end{itemize}

    \item С помощью утилиты VisualVM провести мониторинг и профилирование программы:
    \begin{itemize}
        \item снять график изменения показаний разработанных MBean-классов с течением времени;
        \item определить имя класса, объекты которого занимают наибольший объём памяти JVM, и определить пользовательский класс, в экземплярах которого находятся эти объекты.
    \end{itemize}

    \item С помощью VisualVM и профилировщика IDE локализовать и устранить проблемы с производительностью в программе. В отчёте привести описание выявленной проблемы, способ её устранения и подробный алгоритм локализации со скриншотами. Обеспечить воспроизводимость процесса поиска проблемы.
\end{enumerate}

\newpage
\section*{Реализация}
\addcontentsline{toc}{section}{Реализация}

\subsection*{1. Исходный код MBean}
\addcontentsline{toc}{subsection}{1. Исходный код MBean}

В приложении реализованы два стандартных MBean. \texttt{PointStatistics} хранит число проверенных точек, попаданий и промахов и формирует JMX-уведомление после каждой десятой точки. \texttt{Area} возвращает текущий радиус и площадь заданной вариантом области. Оба MBean используют общее потокобезопасное состояние и регистрируются в платформенном \texttt{MBeanServer} при развёртывании приложения.

\begin{itemize}
    \item \texttt{mbean/PointStatisticsMBean.java}
\end{itemize}
\lstinputlisting[style=java]{src/main/java/com/example/lab3/mbean/PointStatisticsMBean.java}

\begin{itemize}
    \item \texttt{mbean/PointStatistics.java}
\end{itemize}
\lstinputlisting[style=java]{src/main/java/com/example/lab3/mbean/PointStatistics.java}

\begin{itemize}
    \item \texttt{mbean/AreaMBean.java}
\end{itemize}
\lstinputlisting[style=java]{src/main/java/com/example/lab3/mbean/AreaMBean.java}

\begin{itemize}
    \item \texttt{mbean/Area.java}
\end{itemize}
\lstinputlisting[style=java]{src/main/java/com/example/lab3/mbean/Area.java}

\begin{itemize}
    \item \texttt{mbean/MonitoringState.java}
\end{itemize}
\lstinputlisting[style=java]{src/main/java/com/example/lab3/mbean/MonitoringState.java}

\begin{itemize}
    \item \texttt{mbean/MonitoringRegistry.java}
\end{itemize}
\lstinputlisting[style=java]{src/main/java/com/example/lab3/mbean/MonitoringRegistry.java}

\begin{itemize}
    \item \texttt{mbean/MBeanRegistrationListener.java}
\end{itemize}
\lstinputlisting[style=java]{src/main/java/com/example/lab3/mbean/MBeanRegistrationListener.java}

Регистрация новой точки выполняется после сохранения результата в базе данных вызовом \texttt{MonitoringRegistry.registerPoint(result.isHit(), result.getR())}. При изменении радиуса вызывается \texttt{MonitoringRegistry.updateRadius(r)}.

Площадь области вычисляется как сумма площадей прямоугольника, четверти круга и прямоугольного треугольника:
\[
S = \frac{r^2}{2} + \frac{\pi r^2}{16} + \frac{r^2}{4}
  = r^2\left(0{,}75 + \frac{\pi}{16}\right).
\]

\newpage
\subsection*{2. JConsole}
\addcontentsline{toc}{subsection}{2. JConsole}

Приложение развёрнуто на WildFly. Доступ к интерфейсу управления выполнялся через SSH-туннель, а подключение к JMX — по адресу \texttt{service:jmx:remote+http://127.0.0.1:10990}. В дереве MBeans зарегистрированы объекты \texttt{com.example.lab4:type=PointStatistics} и \texttt{com.example.lab4:type=Area}.

\begin{itemize}
    \item Метаданные разработанных MBean.
\end{itemize}
\reportfigure{jconsole-metadata-area.png}{Метаданные MBean Area в JConsole}
\reportfigure{jconsole-metadata-point-statistics.png}{Метаданные MBean PointStatistics в JConsole}

\begin{itemize}
    \item Показания разработанных MBean.
\end{itemize}
\reportfigure{jconsole-attributes-point-statistics.png}{Атрибуты MBean PointStatistics в JConsole}
\reportfigure{jconsole-attributes-area.png}{Атрибуты MBean Area в JConsole}

\begin{itemize}
    \item Время работы виртуальной машины.
\end{itemize}
\reportfigure{jconsole-vm-summary.png}{Сводная информация о виртуальной машине в JConsole}
\reportfigure{jconsole-uptime.png}{Значение Uptime стандартного RuntimeMXBean в миллисекундах}

\begin{itemize}
    \item Получение уведомления.
\end{itemize}
\reportfigure{jconsole-notification.png}{JMX-уведомление при достижении количества точек, кратного 10}

\begin{itemize}
    \item Графики изменения показаний MBean.
\end{itemize}
\reportfigure{jconsole-mbeans-graph-point-statistics.png}{Графики атрибутов PointStatistics в JConsole}
\reportfigure{jconsole-mbeans-graph-area.png}{Графики атрибутов Radius и Area в JConsole}

\begin{itemize}
    \item Выводы.
\end{itemize}

В JConsole проверяется соответствие атрибутов \texttt{TotalPoints}, \texttt{HitPoints} и \texttt{MissPoints} действиям пользователя. Для каждой выборки должно выполняться равенство \texttt{TotalPoints = HitPoints + MissPoints}. При переходе общего счётчика к значению, кратному 10, MBean \texttt{PointStatistics} отправляет уведомление типа \texttt{com.example.lab4.point.totalMultipleOfTen}. Время после запуска JVM определяется по атрибуту \texttt{Uptime} объекта \texttt{java.lang:type=Runtime}; на момент измерения оно составило 19547158 мс.

\newpage
\subsection*{3. VisualVM}
\addcontentsline{toc}{subsection}{3. VisualVM}

VisualVM подключён к той же JVM WildFly через встроенный JMX-интерфейс сервера. Для просмотра пользовательских MBean установлен модуль MBeans.

\begin{itemize}
    \item Обзор JVM.
\end{itemize}
\reportfigure{visualvm-overview.png}{Общая информация о JVM WildFly в VisualVM}

\begin{itemize}
    \item Показания разработанных MBean и их изменение во времени.
\end{itemize}
\reportfigure{visualvm-mbeans-graph-point-statistics-1.png}{Графики изменения атрибутов MBean PointStatistics в VisualVM (часть 1)}
\reportfigure{visualvm-mbeans-graph-point-statistics-2.png}{Графики изменения атрибутов MBean PointStatistics в VisualVM (часть 2)}
\reportfigure{visualvm-mbeans-graph-area.png}{Графики изменения атрибутов Radius и Area в VisualVM}

\begin{itemize}
    \item Классы, занимающие память JVM.
\end{itemize}
\reportfigure{visualvm-memory-classes.png}{Результаты профилирования памяти в VisualVM}
\reportfigure{memory-user-classes.png}{Распределение памяти между классами приложения}

\begin{itemize}
    \item Выводы по результатам мониторинга и профилирования.
\end{itemize}

При добавлении точек значения \texttt{TotalPoints}, \texttt{HitPoints} и \texttt{MissPoints} изменяются в соответствии с результатами проверки. При выборе другого радиуса изменяются атрибуты \texttt{Radius} и \texttt{Area}. В общей гистограмме JVM наибольший объём памяти занимают массивы \texttt{byte[]}. После фильтрации Heap Dump по пакету приложения установлено, что среди пользовательских классов наибольший суммарный shallow size имеет \texttt{CheckResult}: 78 экземпляров занимают 3744 B и удерживают 11064 B. Единственный экземпляр прокси \texttt{AreaCheckBean} удерживает 20832 B.

\newpage
\subsection*{4. Исследование программы на утечки памяти}
\addcontentsline{toc}{subsection}{4. Исследование программы на утечки памяти}

Для воспроизводимого исследования необходимо зафиксировать исходное состояние памяти, создать нагрузку большим количеством проверок точек и повторить измерение после принудительного запуска сборщика мусора. Сравниваются число экземпляров, shallow size и retained size классов, продолжающих удерживаться после GC.

Перед анализом Heap Dump была зафиксирована общая динамика JVM во вкладке \texttt{Monitor}. Пилообразное изменение занятой памяти соответствует циклу выделения объектов и их последующего удаления сборщиком мусора. Такой график показывает активную работу GC, однако сам по себе не определяет объекты, которые продолжают удерживаться, поэтому дальнейшая локализация выполнялась по Heap Dump.

\reportfigure{visualvm-monitor.png}{Мониторинг CPU, Heap, Classes и Threads до локализации проблемы с памятью}

\begin{enumerate}
    \item Подключиться к JVM приложения в VisualVM и открыть вкладку \texttt{Sampler}.
    \item Запустить профилирование памяти и сохранить исходный снимок.
    \item Выполнить серию однотипных действий в приложении, увеличивающих историю результатов.
    \item Нажать \texttt{Perform GC}, повторить снимок и отсортировать классы по занимаемой памяти.
    \item Для растущего класса открыть цепочку удерживающих ссылок и определить пользовательский объект-владелец.
    \item Найти соответствующее поле или коллекцию в исходном коде.
    \item Внести ограничение роста коллекции либо заменить полную загрузку истории постраничной выборкой.
    \item Повторить тот же сценарий нагрузки и сравнить результаты до и после исправления.
\end{enumerate}

До исправления объект \texttt{AreaCheckBean} с областью видимости \texttt{SessionScoped} хранил поле \texttt{List<CheckResult> results}. При инициализации в него загружалась вся история из базы данных, а после каждой проверки новый объект добавлялся в начало списка без удаления старых записей. Heap Dump подтвердил эту причину: поле \texttt{results} ссылалось на \texttt{ArrayList} из 78 элементов с retained size 11544 B.

\reportfigure{memory-before.png}{Использование памяти и удерживаемые объекты до исправления}
\reportfigure{memory-retained-path.png}{Цепочка удержания объектов, позволившая локализовать проблему}
\reportfigure{memory-code-location.png}{Участок исходного кода, вызывающий накопление объектов}
\reportfigure{memory-after.png}{Результат повторного профилирования после исправления}

Для устранения неограниченного роста запрос к базе данных ограничен 50 последними результатами, а после добавления новой записи последний элемент удаляется, если размер списка превысил 50. После повторного развёртывания приложения, аналогичной нагрузки и принудительного запуска сборщика мусора в Heap Dump осталось 50 экземпляров \texttt{CheckResult}. Их shallow size уменьшился с 3744 B до 2400 B, а retained size --- с 11064 B до 7200 B. Retained size объекта \texttt{AreaCheckBean} уменьшился с 20832 B до 11400 B. Таким образом, число удерживаемых результатов и их размер сократились примерно на 36\%, а дальнейший рост списка ограничен.

\newpage
\section*{Вывод}
\addcontentsline{toc}{section}{Вывод}

В ходе лабораторной работы были реализованы и зарегистрированы MBean-классы для наблюдения за количеством проверенных точек, количеством промахов и площадью области. С помощью JConsole проверены значения атрибутов, время работы JVM и получение уведомления после каждой десятой точки. В VisualVM исследованы изменения показателей во времени и распределение памяти JVM. Heap Dump позволил обнаружить неограниченное накопление объектов \texttt{CheckResult} в сессионном объекте \texttt{AreaCheckBean}. Ограничение истории 50 последними результатами устранило дальнейший рост коллекции, что подтверждено повторным профилированием.

\end{document}
